\section{Direct tensor FFT}\label{sec:direct-fft}

This section implements the middle transform of the fractional-shift bridge directly as a complex tensor FFT
$F_{\mathbf t}$ on the $q$-grid. We store the last coordinate as a contiguous record of $r=2^\ell$ complex coefficients
and transform that last axis inside each record. We then run decimation-in-frequency rounds across all $d-1$
record-address axes in parallel, combining the simultaneous butterfly layers of Proposition~\ref{prop:prop14-param} with
the exact positional permutations of Theorem~\ref{thm:positional-perm}.

For Theorem~\ref{thm:main}, computing this
tensor transform directly over $\C$ achieves additive error below $12\ell\cdot 2^{-q}$ with zero output widening. It
needs no Bluestein chirp and no polynomial-ring convolution in~$\Rr_r$, and its error does not grow with the number
$M=R/r$ of records, as it does on the Bluestein route of Appendix~\ref{app:bluestein}. A first-time
reader willing to take the middle transform as a black box needs only Theorem~\ref{thm:direct-fft} for the direct tensor
FFT and Corollary~\ref{cor:walsh-middle} for its bridge interface, and may skip directly to Section~\ref{sec:main}.

\paragraph{Parameters and conventions.}
Throughout this section, let
\[
  t_d=r=2^\ell,\qquad t_i\in\{2^{\ell-1},2^{\ell}\}\ (1\le i<d),\qquad R=t_1\cdots t_d,\qquad M=R/r,\qquad m_0=\lg M .
\]
The last coordinate is stored as the contiguous length-$r$ coefficient list of each record; here $\Rr_r$ specifies only
the record format, and no multiplication in that ring is performed. The exact transform is the normalized positive
tensor DFT $F^+_{\mathbf t}=F_{\mathbf t}$ of Definition~\ref{def:dft}. Stored arrays lie on the $q$-grid $\Cg_q$ with
component width $q+2\le w\le C_wq$. Scalar integer products are delegated to an independent multiplier of cost
$\mathsf M(N)=N\,\mathsf m(N)$ on $N$-bit integers, such as Sch\"onhage--Strassen~\cite{SS71}, and never to a recursive
call of the algorithm of this paper.

\subsection{Phases and scalar products}\label{sec:direct-fft-phases}

This subsection builds the two scalar primitives needed in every FFT round, ensuring that both keep all intermediate
complex numbers inside the closed unit disk.

Because the simultaneous butterfly layer requires disk-array inputs at every round, approximate roots of unity cannot
exceed modulus~$1$ by even a single grid unit. The next lemma constructs a $2r$-entry table of $2r$-th roots of unity on
the $q$-grid by rational series evaluation followed by slight inward shrinking. Each entry has modulus at most~$1$,
error below $2\cdot 2^{-q}$, and deterministic evaluation cost $O(q^3)$ bit operations.

\begin{lemma}[Disk-valued phases]\label{lem:phase-table}
For every precision $q\ge8$ and $r=2^\ell$ with $\ell\le q$, the $2r$ phases $\theta_j=e^{\pi ij/r}$, $0\le j<2r$, have
deterministically computable approximations $\tilde\theta_j\in\Cg_q$ satisfying
\[
  |\tilde\theta_j|\le1,\qquad |\tilde\theta_j-\theta_j|<2\cdot2^{-q}.
\]
The fourth roots of unity are returned exactly. Each entry costs $O(q^3)$ bit operations on dyadic numbers of $O(q)$
bits, on a fixed number of tapes.
\end{lemma}

\emph{Proof idea.} Using trigonometric symmetries, we reduce the angle to $[-\pi/4,\pi/4]$, approximate $\pi$ via
Machin's formula, and sum alternating Taylor series for cosine and sine at $S=2(q+5)+O(\lg q)$ fractional bits.
On non-axis angles, multiplying by $1-2^{-q-3}$ before truncating toward zero keeps the modulus at most~$1$ and the
total error under $2\cdot 2^{-q}$.

\begin{proof}
Put $u_0=2^{-q}$ and $b=q+5$. We compute with dyadic numbers of $S=2b+8\lceil\lg b\rceil+64$ fractional bits, rounding
each product and quotient once with error at most $2^{-S}$, as in part~(iii) of Lemma~\ref{lem:bridge-C} and
Appendix~\ref{app:deferred}. Exact products before a rounding use at most $2S$ fractional bits. We evaluate $\theta_j$
by rational series to within $u_0/16$, then shrink the result slightly and truncate it into the disk.

\emph{Reduction.} Reduce the dyadic rational $\rho=j/r$ without error modulo~$2$ and then, by the symmetries of sine and
cosine, to $|\rho|\le\frac14$. The symmetries contribute only a sign, a conjugation, or a factor $\pm i$, which act
without error; the four fourth roots of unity are recognized directly from the index.

\emph{The angle.} Machin's identity $\pi=16\arctan\frac15-4\arctan\frac1{239}$, summed to $N=4S$ terms with one rounded
division per term, gives $\tilde\pi$ with $|\tilde\pi-\pi|\le21N2^{-S}$, since the omitted series tails lie below
$20\cdot5^{-2N-1}$. Because $\rho$ is dyadic with denominator dividing $r=2^\ell$, the product $x=\tilde\pi\rho$ is
formed without error on the grid of $S+\ell$ fractional bits and rounded once to $S$ fractional bits, adding at most
$2^{-S}$. Thus $|x|<1$ and $|x-\pi\rho|\le(21N+1)2^{-S}$.

\emph{Series.} Sum $\cos x$ and $\sin x$ as the alternating Taylor polynomials $\sum_{n<N}(-1)^nT_n$, with $T_0=1$ and
$T_n=T_{n-1}x^2/((2n)(2n-1))$ for the cosine, and $T_0=x$ and $T_n=T_{n-1}x^2/((2n)(2n+1))$ for the sine. Because the
step multipliers satisfy $x^2/2<\frac13$, every rounding residue is damped, so each sum has rounding error below
$4N2^{-S}$ and tail below $2(1.6\,\mathrm e/N)^N$. Since $\cos$ and $\sin$ are $1$-Lipschitz, the resulting complex
approximation $a$ satisfies
\[
  |a-\theta_j|\le64N2^{-S}+4(1.6\,\mathrm e/N)^N\le2^{-b-4}<\frac{u_0}{16},
\]
as $S-b-4=b+8\lceil\lg b\rceil+60$.

\emph{Into the disk.} Set $\tilde\theta_j=Q_q((1-u_0/8)a)$. Before truncation, the modulus is at most
$(1-u_0/8)(1+u_0/16)<1$ and the error from $\theta_j$ is at most $u_0/16+(u_0/8)(1+u_0/16)<u_0/4$, so
by~\eqref{eq:walsh-truncation} we obtain $|\tilde\theta_j|\le1$ and
$|\tilde\theta_j-\theta_j|<(\sqrt2+\frac14)u_0<2u_0$.

\emph{Cost.} All operands have $O(S)=O(q)$ bits and there are $O(N)=O(q)$ arithmetic steps, so schoolbook arithmetic
takes $O(q^3)$ bit operations per entry without any transcendental oracle.
\end{proof}

Multiplying a disk-grid coefficient by a disk-valued phase must likewise remain inside the unit disk without
intermediate rounding drift. The next lemma shows that forming the complex product without rounding at precision $2q$
and truncating once toward zero preserves $\mathbb D_q$ with error below $\sqrt2\cdot 2^{-q}$.

\begin{lemma}[Disk-preserving products]\label{lem:disk-product}
If $a,b\in\Cg_q$ satisfy $|a|,|b|\le1$, then $Q_q(ab)$ is computed by four exact products of $O(q)$-bit integers and
$O(q)$ further work, lies in the disk $\mathbb D_q$, and satisfies $|Q_q(ab)-ab|<\sqrt2\cdot2^{-q}$.
\end{lemma}

\begin{proof}
Write $a=a_0+ia_1$ and $b=b_0+ib_1$, form $(a_0b_0-a_1b_1)+i(a_0b_1+a_1b_0)$ exactly on the $2q$-bit fractional grid
without rounding the four real products separately, and truncate both components toward zero by $Q_q$. Since $|ab|\le1$,
\eqref{eq:walsh-truncation} gives both claims.
\end{proof}

\subsection{The decimation-in-frequency schedule and its positional maps}\label{sec:direct-fft-layout}

This subsection defines the two positional maps $L$ and~$U$ of the direct tensor FFT and checks that the chunk rounds
between them have the layer shape required by Proposition~\ref{prop:prop14-param}.

The decimation-in-frequency (DIF) recursion runs on all $d-1$ record axes in parallel, one bit level per round. A
positional map $L$ groups their address slots into $K$-bit chunks so that each bulk chunk round becomes a single call to
Proposition~\ref{prop:prop14-param}, and a second positional map $U$ restores natural frequency order afterwards. For a
transform of even length $u$, one positive normalized DIF step applies $H_0$ followed, on the branch $b=1$, by
multiplication by $e^{2\pi ij/u}$:
\begin{equation}\label{eq:dif}
  (F^+_uz)_{2v+b}=\frac{1}{u/2}\sum_{j<u/2}e^{2\pi ivj/(u/2)}\,e^{2\pi ibj/u}\,\frac{z_j+(-1)^bz_{j+u/2}}2,\qquad
  b\in\{0,1\}.
\end{equation}
Successive address slots from high to low thus receive frequency bits from low to high, so the completed DIF recursion
produces bit-reversed frequency order.

\paragraph{Named slots and chunked layout.}
Following the named-slot arrangement of~\cite[Lemma~16]{OAI26}, put $a_i=\lg t_i$ for $1\le i\le d-1$, and name the
address slots of record axis~$i$ from most to least significant by $(i,a_i-1),\dots,(i,0)$. Let
$b_0=(\ell-1)\bmod K$ and $J_{\rm blk}=(\ell-1-b_0)/K$. The original slot list is
$\mathcal O=((i,v):i=1,\dots,d-1,\ v=a_i-1,\dots,0)$. The reordered list begins with the prefix
\[
  \mathcal P=\bigl((i,\ell-1):a_i=\ell,\ i\text{ increasing}\bigr)\ \Vert\
  \bigl((i,v):v=b_0-1,\dots,0\text{ (outer)},\ i\text{ increasing (inner)}\bigr)
\]
and continues with the chunked list
\[
  \mathcal C=\mathop{\Vert}_{j=0}^{J_{\rm blk}-1}\ \mathop{\Vert}_{i=1}^{d-1}
  \bigl((i,v):v=\ell-2-jK,\dots,\ell-1-(j+1)K\bigr),
\]
empty sublists being omitted. Let $L$ be the positional map moving each named slot from its position in $\mathcal O$ to
its position in $\mathcal P\Vert\mathcal C$, and let $B$ be the per-axis bit reversal $(i,v)\mapsto(i,a_i-1-v)$ on
$\mathcal O$.

Once the address slots are reordered into $\mathcal P\Vert\mathcal C$, every bulk DIF round across the $d-1$ record axes
acts on a common bit offset $\rho$ of $d-1$ consecutive $K$-bit chunks, matching the input layout of
Proposition~\ref{prop:prop14-param}. The next lemma verifies this layer shape and shows that applying $U=(LB)^{-1}$
after the rounds untangles both the layout permutation~$L$ and the DIF bit reversal~$B$.

\begin{lemma}[Layout and natural-order restoration]\label{lem:walsh-layout}
In a round at named slot $v$ with $b_0\le v\le\ell-2$, the selected positions in $\mathcal P\Vert\mathcal C$ are bit
$\rho=(v-b_0)\bmod K$, counted from the least significant end, of $D=d-1$ consecutive $K$-bit chunks. With $n_P$ and
$n_S$ denoting the numbers of record-address bits before and after these $d-1$ chunks, the round has layer shape
$[2^{n_P}]\times[2^K]^{d-1}\times[2^{n_S}]$ for Proposition~\ref{prop:prop14-param} (unlike a slice of
Definition~\ref{def:walsh-positional}, whose chunks start at position~$0$, here they start after~$\mathcal P$). The DIF
rounds on the $d-1$ record axes, executed in layout~$L$ on input in the original order, have exact counterpart
$L\,B\,F^+_{\mathbf t_{<d}}$, where $F^+_{\mathbf t_{<d}}$ is the tensor DFT over the first $d-1$ axes. Natural
frequency order is restored by the positional map
\begin{equation}\label{eq:unshuffle}
  U=(LB)^{-1}=B^{-1}L^{-1}.
\end{equation}
Both $L$ and $U$ are compiled in time polynomial in~$q$ and executed exactly by Theorem~\ref{thm:positional-perm}.
\end{lemma}

\emph{Proof idea.} Stripping the top bit of length-$2^\ell$ axes and the lowest $b_0$ bits into $\mathcal P$ leaves an
exact multiple of $K$ bits per record axis. The list $\mathcal C$ interleaves those remaining bits block by block so
that slot~$v$ lands at a common chunk offset~$\rho$, and conjugating the per-axis DIF recursion by $L$ produces
$LBF^+_{\mathbf t_{<d}}$, which $U=(LB)^{-1}$ inverts.

\begin{proof}
\emph{The layout.} The prefix $\mathcal P$ removes the extra top slot of each axis of length $2^\ell$ and the lowest
$b_0$ slots of all $d-1$ record axes. The remaining $\ell-1-b_0=J_{\rm blk}K$ slots of each record axis split into
$J_{\rm blk}$ chunks of width~$K$, which $\mathcal C$ orders first by chunk block $j$ and then by axis $i\in\{1,\dots,d-1\}$.

\emph{The rounds.} In a round at named slot $v$ with $b_0\le v\le\ell-2$, every record axis uses chunk block
$j=\lfloor(\ell-2-v)/K\rfloor$, whose least significant slot is $\ell-1-(j+1)K=(J_{\rm blk}-j-1)K+b_0$. Its offset from
the least significant end of that chunk is therefore $(v-b_0)\bmod K$, while the preceding and following address bits act
as spectators of lengths $n_P$ and~$n_S$.

\emph{The output order.} Applying~\eqref{eq:dif} inductively along each record axis yields the bit-reversed output
order~$B$. Relocating the named slots by $L$ conjugates each physical round by $L$, so executing the rounds in layout
$L$ after applying $L$ to the original-order input computes $LBF^+_{\mathbf t_{<d}}$, proving~\eqref{eq:unshuffle}.

\emph{Compilation.} The slot lists have $m_0=O(q)$ entries, so computing ranks, compositions, and inverses to form $U$ as
$(L\circ B)^{-1}$ takes time polynomial in~$q$.
\end{proof}

The next remark gives a small example showing that the two factors of $U$ in~\eqref{eq:unshuffle} do not commute, and
records that executing $L$ by chunk exchanges, as in~\cite{OAI26}, would reintroduce the cost term that
Theorem~\ref{thm:positional-perm} avoids.

\begin{remark}\label{rem:walsh-unshuffle}
In general $B^{-1}L^{-1}\ne L^{-1}B^{-1}$, and~\eqref{eq:unshuffle} requires the former: for two record axes of bit
widths $(4,3)$ and $K=2$ ($\ell=4$, $b_0=1$), in zero-based destination-list notation $L=(0,3,4,1,5,6,2)$ and
$B=(3,2,1,0,6,5,4)$, so $L\circ B$ sends position~$0$ to~$1$ whereas $B\circ L$ sends it to~$3$. Executing $L$ via
width-$K$ chunk exchanges as in~\cite[Lemma~16]{OAI26} (Appendix~\ref{app:bluestein}) would reintroduce the normalized
cost term $qK^{\tau-1}$; here $L$ and $U$ are executed instead by Theorem~\ref{thm:positional-perm}.
\end{remark}

\subsection{One combined twiddle per record}\label{sec:direct-fft-twiddle}

This subsection combines the up to $d-1$ per-axis twiddle factors of a parallel DIF round into a single root-of-unity
phase per record, so that rounding error does not grow with the dimension~$d$.

Adding their angles exactly, as one integer index into the phase table, before touching the record coefficients ensures
that a round incurs only one phase approximation and one complex multiplication error per coefficient, regardless
of~$d$.

At named slot $v$, each participating record axis $i$ has current subproblem length $u_i=2^{v+1}$. For a record address
$x$, reading the named-slot table gives its current branch bit $b_i=x_{(i,v)}$ and its unprocessed lower index
$j_i^{<v}=\sum_{0\le s<v}2^sx_{(i,s)}$. Because the twiddle of axis~$i$ depends only on the bits of axis~$i$, it
commutes with the butterflies on all other axes, so one parallel round has exact counterpart
\begin{equation}\label{eq:combined-round}
  z\mapsto D_\zeta Hz,\qquad
  \zeta(x)=\exp\Bigl(2\pi i\sum_{i\ \mathrm{participating}}\frac{b_ij_i^{<v}}{u_i}\Bigr),
\end{equation}
where $H$ applies $H_0$ to slot $v$ of each participating axis and $D_\zeta$ multiplies record~$x$ by $\zeta(x)$. Since
every $u_i$ divides $r=2^\ell$, the integer
\[
  J(x)=\sum_{i\ \mathrm{participating}}\frac{r}{u_i}\,b_ij_i^{<v}\bmod r
\]
is computed without error, and the round multiplies every coefficient of record~$x$ by the single phase entry
$\tilde\theta_{2J(x)}$ from Lemma~\ref{lem:phase-table}.

\subsection{The direct tensor FFT}\label{sec:direct-fft-thm}

This subsection assembles the intra-record transform, the layout permutations, and the parallel record-axis rounds into
the complete direct tensor FFT and bounds its error and running time.

Under the parameter regime below, on a tensor~$z$ of the shape fixed above ($t_d=r,\ t_i\in\{r/2,r\}$) with disk-grid
entries in natural tensor order, four sequential stages approximate $F^+_{\mathbf t}z$ directly on the $q$-grid without
widening the component width. First comes an intra-record DIF transform on the last axis, then the layout
permutation~$L$, then at most $\ell$ parallel record-axis butterfly-and-twiddle rounds, and finally the restoring
permutation $U=(LB)^{-1}$. Each of the $\ell$ last-axis levels and at most $\ell$ record-axis rounds adds error below
$6\cdot 2^{-q}$, and the two positional permutations add none, so the total error stays under $12\ell\cdot 2^{-q}$.

\begin{theorem}[Direct complex tensor FFT]\label{thm:direct-fft}
Assume the parameter family of Proposition~\ref{prop:prop14-param} at grid precision~$q$. Suppose $q\ge8$, $d\ge2$,
$r=2^\ell$ with $2\le\ell<q$, $d\ell\le C_\ell q$, $1\le K\le\ell-1$, $t_i\in\{r/2,r\}$ for $1\le i<d$, and
$q+2\le w\le C_wq$. On a disk array on the $q$-grid in natural tensor order, a fixed multitape procedure returns a disk
array on the $q$-grid in natural frequency order on every axis, in component width~$w$ (with no $\lg M$ widening),
satisfying
\begin{equation}\label{eq:direct-fft-error}
  \|\widehat F^+_{\mathbf t}z-F^+_{\mathbf t}z\|_\infty<12\ell\cdot2^{-q},
\end{equation}
in
\begin{equation}\label{eq:direct-fft-cost}
  O\bigl(V[dK+\ell d^{\lambda'}+\ell\,\mathsf M(q+2)/q]\bigr)
  =O\bigl(V[dK+\ell d^{\lambda'}+\ell\,\mathsf m(q+2)]\bigr)
\end{equation}
steps, where $V=2Rw=\Theta(Rq)$.
\end{theorem}

\emph{Proof idea.} The last-axis transform acts inside each length-$r$ record and commutes with record-level operations,
so Lemma~\ref{lem:walsh-layout} identifies the exact composition of the four stages with $F^+_{\mathbf t}$. One
precomputed $2r$-entry phase table supplies all twiddles via sequential tape scans, while bulk record-axis rounds call
Proposition~\ref{prop:prop14-param} and the few remaining prefix rounds apply $H_0$ sequentially at extra precision
before a single truncation. Because every stage keeps entries in the unit disk $\mathbb D_q$ and each exact round
contracts $\|\cdot\|_\infty$, local errors below $6\cdot 2^{-q}$ per level add linearly across $\ell$ last-axis levels
and at most $\ell$ record-axis rounds.

\begin{proof}
The procedure executes four stages sequentially:
\begin{enumerate}[label=(\roman*),leftmargin=*]
  \item Apply a normalized positive DIF FFT on every length-$r$ contiguous record, followed by bit-reversal restoration
    inside each record.

  \item Apply the positional map $L$ on the record addresses via Theorem~\ref{thm:positional-perm}.

  \item Run record-axis rounds at named slots $v=\ell-1,\dots,0$, where at $v=\ell-1$ only axes of length $r$
    participate. Chunk rounds with $b_0\le v\le\ell-2$ call Proposition~\ref{prop:prop14-param} with $D=d-1$, while
    prefix rounds with $v=\ell-1$ or $v<b_0$ apply exact individual butterflies before a single truncation. Each
    completed butterfly round is followed by the combined twiddle multiplication of Section~\ref{sec:direct-fft-twiddle}.

  \item Apply the restoring positional map $U=(LB)^{-1}$ on the record addresses via Theorem~\ref{thm:positional-perm}.
\end{enumerate}
Because the last-axis transform acts within records and commutes with record-axis operations and record permutations,
\eqref{eq:dif} and Lemma~\ref{lem:walsh-layout} show that the exact counterpart of the four stages is
$F^+_{\mathbf t}$.

\emph{Phase table.} One master $2r$-entry phase table from Lemma~\ref{lem:phase-table} serves all rounds. Building it
once costs $rq^{O(1)}=o(V)$, since $M\ge r/2$ and $r$ exceeds every fixed polynomial in~$q$ under the band
$r\ge2^{a_rq/d}$ with $d\le b_dq^\eps$. Every subsequent table lookup is a charged sequential scan, so no random access
is assumed.

\emph{Last axis.} At a DIF level of block length~$u$, stream the two halves of each block onto aligned tapes, form each
half-sum and half-difference without error, and truncate to $\Cg_q$ by $Q_q$. Multiply the $b=1$ branch by its stage
phase via Lemma~\ref{lem:disk-product}. One scan of the master table extracts the stage table
$(\tilde\theta_{(2r/u)j})_{j<u/2}$, which is replayed with a local rewind across the blocks of each record. Splits,
merges, and the final bit reversal on the $\ell$ intra-record bit positions are elementary streams by
Lemma~\ref{lem:transfer-streams}; this intra-record reversal is separate from~$U$. Each level costs
$O(rq+r\,\mathsf M(q+2))$ per record, and the intra-record bit reversal costs $O(rq\ell)$ per record.

\emph{Record axes.} In each chunk round, Lemma~\ref{lem:walsh-layout} gives the layer shape, and setting
$C_M=C_\ell+1$ and $C_{\rm desc}=32(C_\ell+C_w+1)$ satisfies the header hypotheses of
Proposition~\ref{prop:prop14-param}. Indeed, there are $2^{m_0}$ records with $m_0\le d\ell\le C_\ell q$, so
$\lg(2\cdot2^{m_0})\le(C_\ell+1)q$, and the nine-field header has length at most $2(C_\ell+1)q+O(\lg q)\le C_{\rm desc}q$
for large~$q$. In each prefix round, the at most $d-1$ participating $H_0$ factors are applied sequentially without
rounding before one truncation $Q_q$. Intermediate values lie in the disk on denominator exponent at most $q+d-1$, and
temporary sums have modulus at most~$2$, requiring $q+d+O(1)=O(q)$ bits.

In both kinds of record-axis round, for each record~$x$, the index $J(x)$ is computed in time polynomial in~$q$, and
$\tilde\theta_{2J(x)}$ is fetched by one scan of the master table. Every coefficient of the record is then multiplied by
$\tilde\theta_{2J(x)}$ via Lemma~\ref{lem:disk-product}, with total address and scan overhead $O(V)$ by
Lemma~\ref{lem:transfer-long-record}. Every record-axis round thus implements
\begin{equation}\label{eq:round-protocol}
  z\mapsto Q_q\bigl(\tilde D\,Q_q(Hz)\bigr),
\end{equation}
where $\tilde D$ is the recordwise diagonal of stored phases.\footnote{Unlike monomial twists in $\Rr_r$, general
complex twiddles do not commute with componentwise truncation $Q_q$, and the rounded map need not be a contraction
between arbitrary inputs. The error analysis below therefore compares each computed round directly against its exact
contraction $D_\zeta H$ on disk inputs.}

\emph{Error.} Put $u_0=2^{-q}$, let $z$ be a disk array, and set $z'=Q_q(Hz)$. By~\eqref{eq:walsh-truncation},
Lemma~\ref{lem:phase-table}, and Lemma~\ref{lem:disk-product},
\begin{align*}
  \|Q_q(\tilde Dz')-D_\zeta Hz\|_\infty
  &\le\|Q_q(\tilde Dz')-\tilde Dz'\|_\infty+\|(\tilde D-D_\zeta)z'\|_\infty+\|D_\zeta(z'-Hz)\|_\infty\\
  &<(\sqrt2+2+\sqrt2)u_0<6u_0 .
\end{align*}
The same $6u_0$ bound holds for each last-axis level, whose $b=0$ branch incurs one truncation and whose $b=1$ branch
incurs two truncations and one phase error. Every computed intermediate array entering a round is a disk array, because
$H_0$ layers preserve the disk, $Q_q$ does not increase modulus, $|\tilde\theta_j|\le1$, and positional maps permute
values. Since every exact round $A$ is a contraction in $\|\cdot\|_\infty$, for computed input $x$ and ideal input $x_0$
we have $\|\widehat Ax-Ax_0\|_\infty\le\|\widehat Ax-Ax\|_\infty+\|x-x_0\|_\infty$. Summing the local bound $6u_0$ over
$\ell$ last-axis levels and at most $\ell$ record-axis rounds, while $L$ and $U$ act without error by
Theorem~\ref{thm:positional-perm}, gives~\eqref{eq:direct-fft-error}.

\emph{Time.} The at most $b_0+1\le K$ prefix rounds cost $O(VdK)$, and the at most $\ell$ chunk rounds cost
$O(V\ell d^{\lambda'})$. Scalar complex products across the last-axis FFT and the record-axis twiddles cost
$O(V\ell+R\ell\,\mathsf M(q+2))=O(V\ell\,\mathsf M(q+2)/q)$. By Theorem~\ref{thm:positional-perm}, the positional maps
$L$ and $U$ cost $O(V[(m_0/d+K)d^{\lambda'}+K])=O(V\ell d^{\lambda'})$, since $m_0/d\le\ell$ and $K\le\ell-1$. All
subroutines run sequentially on a fixed number of tapes and clean their workspaces on return.
\end{proof}

\subsection{Interface with the bridge}\label{sec:direct-fft-bridge}

This subsection specializes the direct tensor FFT to the grid precision of the fractional-shift bridge and verifies that
it fulfills the middle-transform contract.

Instantiating Theorem~\ref{thm:direct-fft} at the bridge's grid precision $q=P+2$ and component width $w=q+2$, where $P$
is the internal working precision of Theorem~\ref{thm:bridge-precision}, supplies the middle transform used in
Algorithm~\ref{alg:main}. The corollary below records that its error factor is $E_{\rm mid}=12\ell$ and its output widening is
$\delta_{\rm mid}=0$, both well within the limits of Definition~\ref{def:middle-contract}.

\begin{corollary}[Walsh middle transform]\label{cor:walsh-middle}
Assume the hypotheses of Theorem~\ref{thm:direct-fft} at the bridge's grid precision $q=P+2$ with component width
$w=q+2$ (verified for the root family in Lemma~\ref{lem:root-family}, with Lemma~\ref{lem:walsh-lift} covering the
lifted precision $q+8d$ of the internal controlled-flip batches). Applied to a complete tensor $U$ in natural axis order
$(j_1,\dots,j_d)$ with the last axis contiguous and entries in $\mathbb D_q$ (in particular the bridge's prepared spread
array, which satisfies $\|U\|_\infty\le1/8$), Theorem~\ref{thm:direct-fft} returns a disk array on the $q$-grid in
natural frequency order $(u_1,\dots,u_d)$, in component width $q+2$, satisfying
\begin{equation}\label{eq:walsh-middle-error}
  \|\widehat F^+_{\mathbf t}U-F^+_{\mathbf t}U\|_\infty<12\ell\cdot2^{-q}.
\end{equation}
In particular, it satisfies the middle-transform contract (Definition~\ref{def:middle-contract}) of
Theorem~\ref{thm:bridge-precision} with $E_{\rm mid}=12\ell$ and output widening $\delta_{\rm mid}=0$.
\end{corollary}

\begin{proof}
Immediate from Theorem~\ref{thm:direct-fft} and Proposition~\ref{prop:cnot-batch} at grid precision~$q$ and width
$w=q+2$; internal guard allocations and temporary lifts to $q+8d$ return to width $q+2$ at completed boundaries.
\end{proof}

The middle transform uses Theorem~\ref{thm:positional-perm} for $L$ and $U$, where every record carries $r$
coefficients. By contrast, Lemma~\ref{lem:bridge-D} implements the bridge's line-spreading axis exposures by the chunk
interchanges of Lemma~\ref{lem:transfer12} at cost $O(Vd(1+\ell^\tau))$, since a single-line operation has no long
retained record payload.
